\chapter{Topological Spaces}
\pagenumbering{arabic}

\section{Topological Spaces}
\begin{definition}\label{def:topology}
    A \textui{topology} on a set $X$ is a collection $T$ of subsets of $X$ such that:
        \begin{enumerate}[label = (\arabic*),itemsep=1pt,topsep=3pt]
            \item $\emptyset, X \in T$;
            \item if $\{U_i\}_{i \in I} \subseteq T$, then $\bigcup_{i \in I}U_i \in T$;
            \item if $\{U_i\}_{i = 1}^n \subseteq T$, then $\bigcap_{i = 1}^n \in T$.
        \end{enumerate}
    We call the pair $(X,T)$ a \textui{topological space}. The subsets of $X$ contained in $T$ are called \textui{open sets}.
\end{definition}

\begin{example}
    Let $X = \{a,b,c\}$. Then $T_1 := \{\emptyset, \{a,b,c\},\{b\},\{a,b\},\{b,c\}\}$ and $T_2 := \{\emptyset,\{b\},\{a,b,c\}\}$ are topologies on $X$. However, $T_3 =\{\emptyset, \{a\},\{b\},\{a,b,c\}\}$ is \textit{not} a topology on $X.$ Note that $\{a\} \cup \{b\} = \{a,b\}$, which is not in $T_3$.
\end{example}

\begin{example}[Every set admits a topology]
    Let $X$ be a set. The \textit{discrete topology} on $X$ is $\cP(X)$. The \textit{trivial topology} on $X$ is $\{\emptyset, X\}$.
\end{example}

\begin{definition}
    Let $T_1$ and $T_2$ be topologies on a set $X$. If $T_1 \subseteq T_2$, then we say $T_2$ is \textui{finer} than $T_1$, and likewise that $T_1$ is \textui{coarser} than $T_2$.
\end{definition}

\begin{example}\label{example:metric-on-Rn}
    Let $\vec{x} = (x_1,...,x_n)$ and $\vec{y} = (y_1,...,y_n)$ be vectors in $\bfR^n$. Define $d(\vec{x},\vec{y}) := \sqrt{(x_1-y_1)^2 + ... + (x_n - y_n)^2}$. This is a metric on $\bfR^n$, and we will see this induces a topology on $\bfR^n$.
\end{example}

\begin{definition}
    Let $(X,d)$ be a metric space, $x \in X$, and $\epsilon > 0$. The \textui{open ball centered at $x$} \textui{with radius $\epsilon$} is denoted $B(x,\epsilon) = \{y \in X \mid d(x,y) < \epsilon\}$.
\end{definition}

\begin{definition}
    Let $(X,d)$ be a metric space. A set $U \subseteq X$ is \textui{open} if for all $x \in U$, there exists $\epsilon >0$ such that $B(x,\epsilon) \subseteq U$. The collection of open sets of $X$ is denoted $\tau_X$.
\end{definition}

\begin{proposition}
    Let $(X,d)$ be a metric space. Then $(X,\tau_X)$ is a topological space.
\end{proposition}
    \begin{proof}
        Clearly $\emptyset,X \in \tau_X$. Let $\{U_i\}_{i \in I} \subseteq X$ be an arbitrary family of open sets. Let $x \in \bigcup_{i \in i}U_i$. Then $x \in U_i$ for some $i$. So there exists an $\epsilon > 0$ such that $B(x,\epsilon) \subseteq U_i \subseteq \bigcup_{i \in I}U_i$. Now let $\{U_i\}_{i = 1}^n$ be a finite collection of open sets. Let $x \in \bigcap_{i = 1}^n U_i$. Then $x \in U_i$ for all $i = 1,...,n$. So there exists $\epsilon_i > 0$ such that $B(x,\epsilon_i) \subseteq U_i$ for all $i = 1,...,n$. Define $\epsilon := \min\{\epsilon_1,...,\epsilon_n\}$. Then $B(x,\epsilon) \subseteq U_i$ for all $i=1,...,n$. Hence $B(x,\epsilon) \subseteq \bigcap_{i = 1}^n U_i$. Thus $(X,\tau_X)$ is a topological space.
    \end{proof}

\begin{example}
    Using the same terminology from Example~\ref{example:metric-on-Rn}, since $d(\vec{x},\vec{y})$ is a metric, $(\bfR^n, \tau_{\bfR^n})$ is a topological space.
\end{example}

\section{Topologies Generated by a Basis}
\begin{definition}
    Let $X$ be a set. A \textui{basis for a topology on $X$} is a collection $\cB$ of subsets of $X$ such that:
        \begin{enumerate}[label = (\arabic*),itemsep=1pt,topsep=3pt]
            \item for all $x \in X$, there exists $B \in \cB$ such that $x \in \cB$.
            \item If $x \in B_1 \cap B_2$, there exists $B_3 \in \cB$ such that $x \in B_3 \subseteq B_1 \cap B_2$.
        \end{enumerate}
\end{definition}

\begin{remark}
    Let $(X,d)$ be a metric space and $\cB$ a basis for a topology on $X$. It follows from the above definition that for any $U \in \tau_X$, there exists an open basis element $B \in \cB$ such that $x \in B \subseteq U$.
\end{remark}

\begin{example}\label{example:basis-1}
    The sets $\cB_1 := \{(a,b) \mid a<b\}$ and $\cB_2 := \{(a,b) \mid a < b, a,b \in \bfQ\}$ are bases for a topology on $\bfR$. These definitions extend nicely to $\bfR^n$ by considering open balls.
\end{example}

\begin{example}[Every set admits a basis for a topology]\label{example:basis-2}
    Let $X$ be a set. Then $\cB := \{\{p\} \mid p \in X\}$ is a basis for a topology on $X$.
\end{example}

\begin{proposition}\label{prop:top-gen-by-basis}
    Let $X$ be a set and $\cB$ a basis for a topology on $X$. Let $T$ be the set containing the empty set and every subset of $X$ which is equal to a union of basis elements of $\cB$. Then $(X,T)$ is a topological space.
\end{proposition}
    {\color{red} \begin{proof}
        
    \end{proof}}

\begin{definition}
    Let $X$ be a set and $\cB$ a topology on $X$. The topology defined in Proposition~\ref{prop:top-gen-by-basis} is called the \textui{topology generated by $\cB$}.
\end{definition}

\begin{example}
    From Example~\ref{example:basis-1}, $\cB_1$ and $\cB_2$ generate the standard topology on $\bfR^n$. From Example~\ref{example:basis-2}, $\cB$ generates the discrete topology on $X$.
\end{example}

\section{Product Topologies}

\begin{definition}
    Let $(X,T_1)$ and $(Y,T_2)$ be topological spaces. The \textui{product topology on $X \times Y$} is the topology generated by the basis $\cB := \{U \times V \mid U \in T_1, V \in T_2\}$.
\end{definition}

\begin{remark}
    Note that $\cB$ is not necessarily a topology on $X \times Y$. Indeed, 
    consider $X = Y = \{0,1\}$ in the discrete topology. Then the sets $\{0\}\times\{0\}$ and $\{0\}\times\{1\}$ are in $\cB$ as per the above definition, but their union $\{(0,0),(0,1)\}$ is not. Hence $\cB$ is not a topology.
\end{remark}

\begin{example}\label{example:product-bases}
    Recall from Example~\ref{example:basis-1} that:
        \begin{equation*}
        \begin{split}
            \cB := \{(a,b) \mid a < b\}
        \end{split}
        \end{equation*} 
    generates the standard topology on $\bfR$. Consider the following to bases for a topology on $\bfR^2$:
        \begin{equation*}
        \begin{split}
            \cB_1 &:= \{B(x,r) \mid x \in \bfR^2, r > 0\}, \\
            \cB_2 &:= \{I_1 \times I_2 \mid I_1, I_2 \in \cB\}.
        \end{split}
        \end{equation*}
    The basis elements of $\cB_1$ are open balls of radius $r$, whilst the basis elements of $\cB_2$ are open squares. Note that the basis $\cB_2$ interprets $\bfR^2$ as the product $\bfR \times \bfR$.

    This idea can be further generalized to $\bfR^3$. Consider:
        \begin{equation*}
        \begin{split}
            \cB_3 &:= \{B(x,r) \mid x \in \bfR^3 , r > 0\} \\
            \cB_4 &:= \{I_1 \times I_2 \times I_3 \mid I_1,I_2,I_3 \in \cB\} \\
            \cB_5 &:= \{ B \times I \mid B \in \cB_1, I \in \cB\}.
        \end{split}
        \end{equation*}
    The basis $\cB_4$ interprets $\bfR^3$ as $\bfR \times \bfR \times \bfR$, whilst the basis $\cB_5$ interprets $\bfR^3$ as $\bfR^2 \times \bfR$.
\end{example}

\begin{proposition}
    The bases $\cB_1$ and $\cB_2$ from Example~\ref{example:product-bases} generate the same topologies on $\bfR^2$.
\end{proposition}
    {\color{red} \begin{proof}
        
    \end{proof}}

\section{Subspace Topology}

\begin{proposition}\label{prop:subspace-top}
    Let $(X,T)$ be a topological space. If $Y \subseteq X$, then $T_Y := \{Y \cap U \mid U \in T_X\}$ is a topology on $Y$.
\end{proposition}
    {\color{red} \begin{proof}
        
    \end{proof}}

\begin{definition}
    Let $(X,T)$ be a topological space and $Y \subseteq X$. The \textui{subspace topology on $Y$} is the topology defined in Proposition~\ref{prop:subspace-top}. We call $Y$ a \textui{subspace of $X$} if it is paired with the subspace topology.
\end{definition}

\begin{example}
    Consider the interval $[2,4)$ with the subspace topology inherited from $\bfR$. Note that $(3,4)$ is open in $[2,4)$ and $\bfR$, while $[2,3)$ is open in $[2,4)$ but \textit{not} open in $\bfR$.
\end{example}

\begin{example}
    Consider the integers $\bfZ$ with the subspace topology inherited from $\bfR$. For any $n \in \bfZ$, we have $\{n\} = \bfZ \cap \left( n - \frac{1}{2}, n+ \frac{1}{2} \right)$. Since $\left( n - \frac{1}{2}, n+ \frac{1}{2} \right)$ is open in $\bfR$, each singleton in $\bfZ$ is open. Hence the subspace topology on $\bfZ$ inherited from $\bfR$ is the discrete topology.
\end{example}

\begin{proposition}
    If $\cB$ is a basis for a topology on a set $X$, and $Y$ is a subset of $X$, then $\cB_y := \{Y \cap U \mid U \in \cB\}$ is a basis for $Y$.
\end{proposition}
    {\color{red} \begin{proof}
        
    \end{proof}}

\begin{example}
    Recall that $\cB := \{(a,b) \mid a < b\}$ is a basis for the standard topology on $\bfR$. Consider $\cB_\bfQ := \{(a,b) \cap \bfQ \mid a < b\}$. Since $\bfQ$ is dense in $\bfR$, we have that $(a,b) \cap \bfQ \neq \emptyset$. In fact, each nonempty open set of $\bfQ$ with this particular subspace topology is infinite, hence it is not the discrete topology.
\end{example}

\section{Closed Sets}

\begin{definition}
    A subset $A$ of a topological space $X$ is \textui{closed in $X$} if $A^c$ is open.
\end{definition}

\begin{example}
    Let $(X,T)$ be a topological space. Then $\emptyset$ and $X$ are closed.
\end{example}

\begin{example}\label{example:singletons-closed-in-Rn}
    If $x \in \bfR^n$, then $\{x\}$ is closed in $\bfR^n$. Indeed, for any $p \in \bfR^n \setminus \{x\}$, we have that $p \in B \left( p , \frac{\lnorm p - x \rnorm}{3} \right) \subseteq \bfR^n\setminus\{x\}$. Since $\bfR^n \setminus \{x\}$ is open, $\{x\}$ must be closed.
\end{example}

\begin{example}
    Given any $x \in \bfR^n$, the set $\{y \in \bfR^n \mid \lnorm x - y \rnorm \leq \epsilon\}$ is closed. The proof of this fact follows identically to Example~\ref{example:singletons-closed-in-Rn}.
\end{example}

\begin{example}
    Consider the topological space $(X,\cP(X))$. Since every set is open, every set must be closed.
\end{example}

\begin{proposition}
    Let $X$ be a topological space. Then:
    \begin{enumerate}[label = (\arabic*),itemsep=1pt,topsep=3pt]
        \item $\emptyset$ and $X$ are closed;
        \item If $\{C_i\}_{i \in I}$ is an arbitrary family of closed sets, then $\cap_{i \in I}C_i$ is closed;
        \item If $\{C_i\}_{i =1}^n$ is a finite collection of closed sets, then $\bigcup_{i = 1}^n C_i$ is closed.
    \end{enumerate}
\end{proposition}
    \begin{proof}
        This follows from De Morgan's laws and Definition~\ref{def:topology}.
    \end{proof}

\begin{example}
    Every finite subset of $\bfR^n$ is closed. Indeed, any finite subset can be written as the finite union of singletons, which are closed in $\bfR^n$ by Example~\ref{example:singletons-closed-in-Rn}.
\end{example}

\begin{proposition}
    Let $Y$ be a subspace of a topological space $X$. Then a set $A \subseteq Y$ is closed in $Y$ if and only if $A = Y \cap C$, where $C$ is a closed set in $X$.
\end{proposition}
    {\color{red} \begin{proof}
        
    \end{proof}}

\section{Closure and Interior}

\begin{definition}
    Let $(X,T)$ be a topological space and $A\subseteq X$.
    \begin{enumerate}[label = (\arabic*),itemsep=1pt,topsep=3pt]
        \item The \textui{interior of $A$} is $A^o := \bigcup \{U \mid U \subseteq A, U \text{ open }\}$.
        \item The \textui{closure of $A$} is $\overline{A} := \bigcap\{C \mid C \supseteq A, C \text{ closed }\}$.
    \end{enumerate}
\end{definition}

\begin{remark}
Clearly $A^o \subseteq A \subseteq \overline{A}$. The interior of any set is open and is the largest open subset containing it. The closure of any set is closed and is the smallest closed subset contained in it. A set is open if and only if it is equal to its interior. A set is closed if and only if it is equal to its closure. If $C$ is any closed set in $X$ and $A \subseteq C$, then $\overline{A} \subseteq C$.
\end{remark}

\begin{example}
    If $F \subseteq \bfR^n$ is finite then $F^o = \emptyset$ and $\overline{F} = F$.
\end{example}

\begin{example}
    Consider the interval $(2,3]$ in $\bfR$. Then $(2,3]^o = (2,3)$ and $\overline{(2,3]} = [2,3]$.
\end{example}

\begin{example}
    In the discrete topology, every set is open and closed, hence every set equals its interior and closure.
\end{example}

\begin{example}
    Consider the set $A:= \{1,\frac{1}{2}, \frac{1}{3},...\}$ as a subset of $\bfR$. Then $A^o = \emptyset$ and $\overline{A} = A \cup \{0\}$.
\end{example}

\begin{proposition}
    Let $(X,T)$ be a topological space and $A \subseteq X$.
    \phantom{a}
    \begin{enumerate}[label = (\arabic*),itemsep=1pt,topsep=3pt]
        \item $p \in \overline{A}$ if and only if for all open sets $U$ in $X$ containing $p$, we have $U \cap A \neq \emptyset$.
        \item If $\cB$ is a basis for $T$, then $p \in \overline{A}$ if and only if for all basis elements $B \in \cB$ containing $p$, we have $B \cap A \neq \emptyset$.
    \end{enumerate}
\end{proposition}
    {\color{red} \begin{proof}
        
    \end{proof}}

\section{Limit Points of a Set}

\begin{definition}
    Let $(X,T)$ be a topological space. Let $A \subseteq X$. A point $p \in X$ is a \textui{limit point of $A$} if $p \in \overline{A \setminus \{p\}}$. The set of limit points of $A$ is denoted $A'$.
\end{definition}

\begin{example}
    Consider $\bfR$ with the standard topology.
    \begin{enumerate}[label = (\arabic*),itemsep=1pt,topsep=3pt]
        \item If $F$ is a finite set, then $F' = \emptyset$.
        \item For any $a<b$, we have $(a,b]' = [a,b]$.
        \item 
        \item The set containing the limit points of $\{\frac{1}{n} \mid n \in \bfN\}$ is $\{0\}$.
    \end{enumerate}
\end{example}

\begin{example}
    Let $X$ be a set with the discrete topology. For any $A \subseteq X$, the set containing the limit points of $A$ is empty.
\end{example}

\begin{proposition}\label{prop:closed-iff-contains-lim-pts}
    Let $(X,T)$ be a topological space. Let $A \subseteq X$. Then $A$ is closed in $X$ if and only if $A$ contains all of its limit points.
\end{proposition}
    \begin{proof}
        \begin{equation*}
        \begin{split}
            \text{A is closed in $X$}
            &\iff \text{$A^c$ is open in $X$} \\
            &\iff \text{$(\forall p \in A^c)(\exists U \in T):p \in U \subseteq A^c$} \\
            &\iff \text{$(\forall p \in A^c)(\exists U \in T):p \in U $ and $ U \cap A= \emptyset$} \\
            &\iff \text{$(\forall p \in A^c):p \not\in \overline{A}$} \\
            &\iff \text{$(\forall p \in A^c):p \not\in \overline{A \setminus \{p\}}$} \\
            &\iff \text{$(\forall p \in A^c):p \not\in \overline{A \setminus \{p\}}$} \\
            &\iff \text{$A$ contains its limit points.} \qedhere 
        \end{split}
        \end{equation*}
    \end{proof}

\begin{proposition}
    Let $(X,T)$ be a topological space and $A \subseteq X$. Then $\overline{A} = A \cup A'$.
\end{proposition}
    \begin{proof}
        Let $p \in \overline{A}$. If $p \in A$, then by Proposition~\ref{prop:closed-iff-contains-lim-pts} we are done. Suppose $p \in \overline{A} \setminus A$. Since $p \in \overline{A}$, for all open sets $U \subseteq X$ containing $p$, we have $U \cap A \neq \emptyset$. But since $p \not\in A$, this is equivalent to the statement that for all open sets $U \subseteq X$ containing $p$, we have $U \cap (A\setminus\{p\}) \neq \emptyset$. Therefore $p \in \overline{A \setminus \{p\}}$. This establishes $\overline{A} \subseteq A \cup  A'$.

        Now let $p \in A \cup A'$. If $p \in A$, then $p \in \overline{A}$. If $p \in A'$, then $p \in \overline{A}$ by Proposition~\ref{prop:closed-iff-contains-lim-pts}. Hence $\overline{A} = A \cup A'$.
    \end{proof}

\section{Hausdorff Spaces}

\begin{definition}
    Let $(X,T)$ be a topological space. We say $X$ is \textui{Hausdorff} (or \textui{$T_2$}) if for all $x,y \in X$, there exist an open set $U \subseteq X$ containing $x$ and an open set $V \subseteq X$ containing $Y$ such that $U \cap V = \emptyset$.
\end{definition}

\begin{example}
    The space $(\bfR^n, \tau_{\bfR^n})$ is Hausdorff. For any $x,y \in \bfR^n$, let $\epsilon := \frac{\lnorm x-y \rnorm}{3}$. Then $B(x,\epsilon) \cap B(y,\epsilon) = \emptyset$.
\end{example}

\begin{example}
    The space $(X,\cP(X))$ is Hausdorff since for any $x,y \in X$, we have $\{x\} \cap \{y\} = \emptyset$.
\end{example}

\begin{example}
    The trivial topology on a set $X$ is not Hausdorff.
\end{example}

\begin{example}
    Let $X := \{a,b,c\}$ and $T:= \{\emptyset, \{b\}, \{a,b\}, \{b,c\}, \{a,b,c\}\}$. Then $(X,T)$ is not Hausdorff. Indeed, the elements $b$ and $c$ cannot be separated by disjoint open sets, as any open set which contains $c$ also contains $b$.
\end{example}

\begin{proposition}
    In any Hausdorff space, every finite set is closed.
\end{proposition}
    \begin{proof}
        Let $p,q \in X$. Then there exists an open set $U$ containing $p$ and an open set $V$ containing $q$ such that $U \cap V = \emptyset$. Now $p \in X \setminus \{q\}$. Hence $p \in U \subseteq X \setminus \{q\}$. Therefore $X \setminus\{q\}$ is open, establishing that $\{q\}$ is closed. Since any finite set can be written as the finite union of singletons, every finite set in a Hausdorff space is closed.
    \end{proof}

\begin{proposition}
    Let $(X,T)$ be Hausdorff and $A \subseteq X$. Then a point $p \in X$ is a limit point of $A$ if and only if every open set $U \subseteq X$ which contains $p$ also contains infinitely many points of $A$.
\end{proposition}
    \begin{proof}
        We will first prove the backwards direction. Let $U \subseteq X$ be open and $p \in U$. If $U$ contains infinitely many points of $A$, then $U$ will contain infinitely many points of $A \setminus \{p\}$. Hence $U \cap (A \setminus \{p\}) \neq \emptyset$. Since $U$ was arbitrary, $p \in \overline{A \setminus \{p\}}$, which means $p$ is a limit point of $A$.

        Now let $p$ be a limit points of $A$. Suppose towards contradiction there exists an open set $U \subseteq X$ containing $p$ such that $\card(U \cap A) < \infty$. Then we also have $\card\bigl(U \cap (A \setminus \{p\})\bigr) < \infty$. Suppose $U \cap (A \setminus \{p\}) = \{a_1,...,a_n\}$. Since $X$ is Hausdorff, $X \setminus \{a_1,...,a_n\}$ is open. In particular, $U \cap (X \setminus \{a_1,...,a_n\})$ is open and contains $p$. But this means:
            \begin{equation*}
            \begin{split}
                \Bigl(U \cap (X \setminus\{a_1,...,a_n\})\Bigr) \cap (A \setminus \{p\}) 
                & = \Bigl(U \cap (A \setminus \{p\})\Bigr) \cap (X \setminus\{a_1,...,a_n\}) \\
                & = \{a_1,...,a_n\} \cap (X \setminus\{a_1,...,a_n\}) \\
                & = \emptyset.
            \end{split}
            \end{equation*}
        Since there exists an open set containing $p$ which is disjoint from $A \setminus \{p\}$, $p \not\in \overline{A \setminus \{p\}}$. This is a contradiction, as we assumed $p$ to be a limit point of $A$.
    \end{proof}

\begin{definition}
    In an arbitrary topological space $X$, a sequence $\{x_n\}_{n = 1}^\infty \subset X$ \textui{converges} to a point $x \in X$ provided that, for all neighborhoods $U$ of $x$, there is a positive integer $N$ such that $x_n \in U$ for all $n \geq N$.
\end{definition}

\begin{theorem}
    If $X$ is a Hausdorff space, then a sequence of points of $X$ converges to at most one point of $X$.
\end{theorem}
    {\color{red} \begin{proof}
        
    \end{proof}}

\begin{example}
    Consider the topological space $(\bfN,T)$, where $T$ is the cofinite topology. Consider the sequence $\{i\}_{i = 1}^\infty \subset \bfN$. Let $p \in \bfN$ be arbitrary. Let $U$ be a neighborhood of $p$. Since $U$ is open in the cofinite topology, $\card(U^c) < \infty$. Pick the largest $N \in U^c$. Then for all $i \geq N$ we have $i \in U$. Therefore the sequence $\{i\}_{i = 1}^\infty$ converges to $p$. Since $p$ was arbitrary, the sequence $\{i\}_{i = 1}^\infty$ converges to every point in $\bfN$.
\end{example}

\section{Continuity}

\begin{definition}
    Let $(X,T_X)$ and $(Y,T_Y)$ be topological spaces. A function $f:X \rightarrow Y$ is \textui{continuous} if for all $U \in T_Y$, we have $f^{-1}(U) \in T_X$.
\end{definition}

\begin{example}
    For topological spaces $(X,\cP(X))$ and $(Y,T)$, any map $f:X \rightarrow Y$ is continuous.
\end{example}

\begin{example}
    Consider the sets $X := \{a,b,c,d\}$ and $Y := \{1,2,3\}$. Let:
        \begin{equation*}
        \begin{split}
            T_X &:= \{\emptyset, \{c\}, \{a,b\}, \{a,b,c\}, \{c,d\} \{a,b,c,d\}\},\\
            T_Y &:= \{\emptyset, \{1\}, \{2\}, \{1,2\}, \{1,2,3\}\}.
        \end{split}
        \end{equation*}
    Then $f:X \rightarrow Y$ given by:
        \begin{equation*}
        \begin{split}
            f(a) &= 1, \\
            f(b) &= 1, \\
            f(c) &= 2, \\
            f(d) &= 2, \\
        \end{split}
        \end{equation*}
    is continuous, while $g:X \rightarrow Y$ given by:
        \begin{equation*}
        \begin{split}
            g(a) &= 1, \\
            g(b) &= 2, \\
            g(c) &= 2, \\
            g(d) &= 3, \\
        \end{split}
        \end{equation*}
    is not continuous.
\end{example}

\begin{proposition}\label{prop:continuity-equivalences}
    Let $X,Y$ be topological spaces. Let $f:X \rightarrow Y$ be a map. The following are equivalent:
        \begin{enumerate}[label = (\arabic*),itemsep=1pt,topsep=3pt]
            \item $f$ is continuous;
            \item For all $A \subseteq X$, $f\bigl(\h1\overline{A}\h1\bigr) \subseteq \overline{f(A)}$;
            \item $f^{-1}(C)$ is closed in $X$ for any $C$ closed in $Y$.
            \item For all $x_0 \in X$ and for all $V \in T_Y$ such that $f(x_0) \in V$, there exists a $U \in T_X$ such that $x_0 \in U$ and $f(U) \subseteq V$.
        \end{enumerate}
\end{proposition}
    {\color{red} \begin{proof}
        
    \end{proof}}

\begin{example}
    The following is a counter example of (2) from Proposition~\ref{prop:continuity-equivalences} when $f$ is not continuous. Define $f:\bfR \rightarrow \bfR$ by $f(x) := \chi_{\bfQ}(x)$. Then $f(\overline{\bfQ}) = f(\bfR) = \{0,1\}$, whilst $\overline{f(\bfQ)} = \overline{\{1\}} = \{1\}$.
\end{example}

\section{Metric Topologies}

\begin{definition}
    A \textui{metric topology} is the topology generated by the basis of open balls defined by the given metric.
\end{definition}

\begin{definition}
    Let $(X,T)$ be a topological space. We say $X$ is \textui{metrizable} if there exists a metric $d$ that induces the given topology on $X$.
\end{definition}

\begin{example}
    The topological space $(X,\cP(X))$ is metrizable.
    \iffalse
     Define $d:X \times X \rightarrow \bfR$ by:
        \begin{equation*}
        \begin{split}
            d(x,y) = \begin{cases} 1, & x\neq y \\ 0, & x = y\end{cases}.
        \end{split}
        \end{equation*}
    This is called the \textit{discrete metric}. Let $T_d$ be the metric topology induced by the metric $d$. Note that 
    \fi
\end{example}


\iffalse
\begin{example}
    Consider the discrete topology $(X,\cP(X))$. Then $X$ is metrizable. Indeed, define $d:X \times X \rightarrow \bfR$ by:
        \begin{equation*}
        \begin{split}
            d(x,y) = \begin{cases} 1, & x\neq y \\ 0, & x = y\end{cases}.
        \end{split}
        \end{equation*}
    This is a metric, and is called the \textit{discrete metric}. Let $T_d$ denote the topology generated by the open balls with respect to the metric $d$. Let $Y \subseteq X$ be open. Then $Y = \bigcup_{y \in Y} \{y\}$. But note that we can write $\{y\} = B_d(y, \frac{1}{2})$. Hence $Y = \bigcup_{y \in Y}B_d(y,\frac{1}{2})$. This gives the inclusion $\cP(X) \subseteq T_d$. Conversely, since every open ball is a subset of $X$, surely it is open in the discrete topology. Hence $T_d \subseteq \cP(X)$.
\end{example}
\fi

\begin{proposition}
    Every metrizable topological space is Hausdorff.
\end{proposition}
    {\color{red} \begin{proof}
        
    \end{proof}}

\begin{example}
    Consider the three point set $\{a,b,c\}$ with the topology 
        \begin{equation*}
        \begin{split}
            T:=\{\emptyset, \{b\}, \{a,b\}, \{b,c\}, \{a,b,c\}\}.
        \end{split}
        \end{equation*}
    Since this topological space is not Hausdorff, it is not metrizable.
\end{example}

\begin{definition}
    Let $(X,d)$ be a metric space and $E \subseteq X$. The \textui{diameter of $E$} is $\diam(E) := \sup\{d(x,y) \mid x,y \in E\}$. If $\diam(E) < \infty$, we say $E$ is \textui{bounded}.
\end{definition}

\begin{proposition}
    Let $(X,d)$ be a metric space. Define $\overline{d}:X \times X \rightarrow \bfR$ by $\overline{d}(x,y):= \min\{d(x,y),1\}$. Then $\overline{d}$ is a metric and induces the same topology as $d$.
\end{proposition}

\begin{theorem}
    Let $X$ be a topological space. Let $A \subseteq X$. If $\{x_n\}_{n = 1}^\infty \subset A$ is a sequence which converges to some $x_0 \in X$, then $x_0 \in \overline{A}$. The converse holds if $X$ is metrizable.
\end{theorem}
    \begin{proof}
        Let $U$ be a neighborhood of $x_0$. Find $N$ sufficiently large such that for all $n \geq N$ we have $x_n \in U$. Since $\{x_n\}_{n = 1}^\infty \subset U$, we have $x_n \in U \cap A$. Since $U \cap A$ is nonempty for every neighborhood of $x_0$, this means $x_0 \in \overline{A}$.

        Conversely, suppose $X$ is metrizable. There exists a metric $d$ which induces the same topology on $X$. As $x_0 \in \overline{A}$, we can inductively find a sequence $\{x_n\}_{n = 1}^\infty$ such that $x_n \in B(x_0, \frac{1}{n}) \cap A$. Hence $\{x_n\}_{n = 1}^\infty \subset A$ and also converges to $x_0$.
    \end{proof}

\begin{example}
    Consider the three point set $X:=\{a,b,c\}$ with the topology 
        \begin{equation*}
        \begin{split}
            T:=\{\emptyset, \{b\}, \{a,b\}, \{b,c\}, \{a,b,c\}\}.
        \end{split}
        \end{equation*}
    If $A \subseteq X$ is a neighborhood of $b$, then $\overline{A} = X$. Hence for each $p \in \overline{A}$, we have a sequence $\{b\}_{n = 1}^\infty$ converging to $p$. Since $X$ is not Hausdorff, it is not metrizable, so the converse of Theorem~\ref{} does not hold.
\end{example}

\section{Quotient Topologies}

\begin{definition}
    Let $X,Y$ be topological spaces and $p:X \rightarrow Y$ a continuous surjection. We say $p$ is a \textui{quotient map} if $U$ is open in $Y$ if and only if $p^{-1}(U)$ is open in $X$.
\end{definition}

\begin{remark}
An equivalent condition is to require that a subset $A$ of $Y$ be closed in $Y$ if and only if $p(A)$ is closed in $X$. It follows from the equation $p^{-1}(Y \setminus B) = X \setminus p^{-1}(B)$.    
\end{remark}

\begin{proposition}
    An open (or closed) continuous surjection is a quotient map.
\end{proposition}
{\color{red} \begin{proof}
    
\end{proof}}

\begin{remark}
Note that continuous surjections need not be quotient maps, and quotient maps need not be open (or closed) maps.
\end{remark}

\begin{example}
    Define $f:(\bfR, \cP(\bfR)) \rightarrow (\bfR,\tau_{\bfR})$ by $f(x) = x$. This map is continuous, surjective, but is not a quotient map. Note that $p^{-1}(\{3\})$ is open in $(\bfR,\cP(\bfR))$ but not in $(\bfR,\tau_{\bfR})$.
\end{example}

\begin{example}
    Let $X:= [0,1] \cup [2,3] \subseteq \bfR$ and $Y = [0,2] \subseteq \bfR$. Define $p:X \rightarrow Y$ by:
        \begin{equation*}
        \begin{split}
            p(x):=\begin{cases}
            x, &x \in [0,1] \\
            x-1, &x \in [2,3]
            \end{cases}.
        \end{split}
        \end{equation*}
    Then $p$ is closed, continuous, and a surjection. Hence it is a quotient map. But $[0,1]$ is open in $X$, whilst $p([0,1]) = [0,1]$ is not open in $Y$. Thus $p$ is not an open map.
\end{example}

\begin{proposition}\label{prop:unique-quotient-top}
    Let $X$ be a topological space and $Y$ any set. Let $p:X \rightarrow Y$ be a surjection. The set $\{U \subseteq Y \mid p^{-1}(U) \text{ open in $X$} \}$ is the unique topology on $Y$ which makes $p$ a quotient map.
\end{proposition}
    {\color{red} \begin{proof}
        
    \end{proof}}

\begin{definition}
    The topology in Proposition~\ref{prop:unique-quotient-top} is called the \textui{quotient topology on $Y$} \textui{induced by $p$}.
\end{definition}

\begin{proposition}
    Let $(X,T)$ be a topological space and $Y$ a set. Let $p:X \rightarrow Y$ be a surjection. The quotient topology on $Y$ induced by $p$ is the finest topology on $Y$ such that $p$ is continuous.
\end{proposition}
    {\color{red} \begin{proof}
        
    \end{proof}}

\begin{example}
    Define $p:\bfR \rightarrow \{a,b,c\}$ by:
        \begin{equation*}
        \begin{split}
            p(x):=\begin{cases}
            a,& x>0 \\
            b,& x < 0 \\
            c,& x = 0
            \end{cases}.
        \end{split}
        \end{equation*}
    Then the set:
        \begin{equation*}
        \begin{split}
            \{X \subseteq \{a,b,c\} \mid p^{-1}(X) \text{ open in $\bfR$ }\}=\{\emptyset, \{a\}, \{b\}, \{a,b\},\{a,b,c\}\}
        \end{split}
        \end{equation*}
    is the quotient topology on $Y$ induced by $p$.
\end{example}

\begin{definition}
    Let $\sim$ be an equivalence relation on a topological space $X$. Let $X /\sim$ denote the set of distinct equivalence classes. Let $p:X \rightarrow X/\sim$ be the canonical surjection given by $p(x) = [x]$. Let $T_\sim$ be the quotient topology on $X /\sim$ induced by $p$. We call the space $(X/\sim,T_\sim)$ the \textui{quotient space of $X$ modulo $\sim$}.
\end{definition}

\begin{example}
    Consider the set $S := \{(x,y) \mid x \in [0,1], y \in [0,1]\}$. Define $\sim$ on $S$ by:
        \begin{equation*}
        \begin{split}
            &(0,y) \sim (1,y) \text{ for all } y \in [0,1] \\
            &(x,y) \sim (x,y) \text{ for all }(x,y) \in (0,1) \times [0,1].
        \end{split}
        \end{equation*}
        This is an equivalence relation. The space $(S/\sim, T_\sim)$ is homeomorphic to a cylinder in $\bfR^3$.
\end{example}

\begin{example}
    Let $D = \{(x,y) \mid x^2 + y^2 \leq 1\}$. Define $\sim$ on $D$ by:
        \begin{equation*}
        \begin{split}
            &(x,y) \sim (x,y) \text{ whenever }x^2 + y^2 < 1,\\
            &(x,y) \sim (a,b) \text{ whenever }x^2 + y^2 = 1.\\
        \end{split}
        \end{equation*}
    This is an equivalence relation. The space $(D/\sim, T_\sim)$ is homeomorphic to a sphere in $\bfR^3$.
\end{example}

\begin{corollary}
    Let $g:X \rightarrow Z$ be a surjective continuous map. Let $X^\ast = \{g^{-1}(\{z\}) \mid z \in Z\}$. Give $X^\ast$ the quotient topology. Then $g$ induces a bijective continuous map $f:X^\ast \rightarrow Z$ which is a homeomorphism if and only if $g$ is a quotient map. Moreover, the following diagram commutes:
    \begin{center}
    \begin{tikzcd}
    X \arrow[d, "p"'] \arrow[r, "g"] & Z \\
    X^\ast \arrow[ru, "f"', dotted]  &  
    \end{tikzcd}
    \end{center}

\end{corollary}
    {\color{red} \begin{proof}
        
    \end{proof}}

\begin{example}
    Define $\sim$ on $[0,1] \times [0,1]$ by:
        \begin{equation*}
        \begin{split}
            (0,y) &\sim (1,y) \text{ for all }y \in [0,1]\\
            (x,0) &\sim (x,1) \text{ for all }x \in [0,1]\\
            (x,y) &\sim (x,y) \text{ for all }(x,y) \in (0,1) \times (0,1).\\
        \end{split}
        \end{equation*}
    We will prove that $([0,1] \times [0,1])/\sim$ is homeomorphic to $S^1 \times S^1$. Define $g:[0,1] \times [0,1] \rightarrow S^1 \times S^1$ by $g(s,t) = (\cos 2\pi s, \sin 2\pi s, \cos 2 \pi t, \sin 2 \pi t)$. This map is continuous and surjective.
\end{example}

\begin{example}
    Let $p:X \rightarrow Y$ be a quotient map and $A \subseteq X$ a subspace. Then $\restr{p}{A} :A \rightarrow p(A)$ is not necessarily a quotient map. Indeed, let $X:=[0,1] \cup [2,3] \subseteq \bfR$ and $Y := [0,2] \subseteq \bfR$. Define $p:X \rightarrow Y$ by:
        \begin{equation*}
        \begin{split}
            p(x):=\begin{cases}
            x, & x \in [0,1] \\
            x-1, & x \in [2,3]
            \end{cases}.
        \end{split}
        \end{equation*}
    Since $p$ is a closed map, it is a quotient map. Consider the subset $A := [0,1) \cup [2,3] \subseteq X$. We have $(\restr{p}{A})^{-1}([1,2]) = [2,3]$. Note that $[2,3]$ is open in $A$ but $[1,2]$ is not open in $Y$.
\end{example}

\begin{example}
    Let $p:A \rightarrow B$ and $q:C \rightarrow D$ be quotient maps. Then the product $p \times q:A \times C \rightarrow B \times D$ is not necessarily a quotient map.
\end{example}
